\chapter{Département d'attache et mission}

\section{Introduction}

Ce chapitre présente le département d’attache au sein de MEDIASSIVE ainsi que les missions qui m’ont été confiées durant mon Projet de Fin d’Études. Il met l’accent sur mon rôle en tant que chef d’équipe stagiaire, développeur full-stack et responsable du déploiement, ainsi que sur ma participation à plusieurs projets internes de l’entreprise, notamment FinderrLink, ScanVivo et Batiscan.

\section{Aperçu sur le département d'attache}

Le département développement web et digital de MEDIASSIVE constitue le pôle technique chargé de concevoir, développer, déployer et maintenir des solutions numériques. Il intervient aussi bien sur des projets clients que sur des projets internes de type SaaS, dans une logique d’innovation et de création de valeur pour l’entreprise.

Le travail au sein de ce département suit généralement plusieurs étapes : analyse des besoins, conception de l’architecture, développement front-end et back-end, intégration de la base de données, tests, correction des anomalies, déploiement et suivi technique. Ces étapes nécessitent une collaboration entre plusieurs profils, notamment les développeurs full-stack, les développeurs IA, les designers UI/UX et les responsables de projet.

\section{Position et rôle dans l'entreprise}

Durant mon stage, j’ai occupé le rôle de chef d’équipe stagiaire au sein du département digital. Mon poste ne se limitait pas uniquement au développement technique, mais comprenait également la coordination avec les autres stagiaires, le suivi de l’avancement des tâches, l’organisation du travail et la vérification de la cohérence entre les différentes parties des projets.

Sur le projet FinderrLink, j’ai principalement travaillé en tant que développeur full-stack et responsable du déploiement. Mes missions ont porté sur la mise en place de la base technique de la plateforme, le développement des premières fonctionnalités, la structuration de l’application, la gestion des rôles et la préparation de l’environnement de déploiement. Les détails fonctionnels de FinderrLink seront présentés dans les chapitres suivants.

L’équipe principale de FinderrLink était composée de plusieurs profils complémentaires. J’ai assuré la partie développement full-stack et déploiement, Mme Chaimae Mezouri a travaillé sur la partie intelligence artificielle, et M. Noeam El Afia a participé à la conception UI/UX de la plateforme. Cette organisation a permis de répartir les tâches selon les compétences de chaque membre.

\section{Organisation des projets et coordination}

En plus de FinderrLink, j’ai également participé aux projets ScanVivo et Batiscan. Dans un premier temps, j’ai développé et déployé les premières versions basiques de ces deux plateformes afin de mettre en place une base fonctionnelle exploitable. Par la suite, ces projets ont été confiés à une équipe de stagiaires à distance, chargée de poursuivre leur développement sous suivi et orientation.

Une autre équipe de stagiaires à distance a été dédiée au développement d’un panneau d’administration pour FinderrLink. Ce panneau avait pour objectif de permettre la supervision de l’application, notamment à travers le suivi des utilisateurs, des missions et des données principales de la plateforme.

Mon rôle consistait à coordonner ces équipes, suivre l’avancement des tâches, expliquer les besoins, orienter les choix techniques et vérifier que le travail réalisé restait cohérent avec les objectifs des projets. Cette organisation m’a permis de développer des compétences en gestion d’équipe, en communication et en suivi de projet, en plus de mes compétences techniques.

\section{Tâches réalisées}

Les principales tâches réalisées durant cette mission peuvent être résumées comme suit :

\begin{itemize}
\item Développement full-stack de la première version de FinderrLink.
\item Préparation et réalisation du déploiement de la plateforme.
\item Mise en place des bases techniques des projets ScanVivo et Batiscan.
\item Déploiement des premières versions basiques de ScanVivo et Batiscan.
\item Coordination avec l’équipe IA et l’équipe UI/UX.
\item Suivi de deux équipes de stagiaires à distance.
\item Accompagnement des stagiaires dans la compréhension des besoins.
\item Suivi du développement du panneau d’administration de FinderrLink.
\item Vérification de la cohérence technique et fonctionnelle entre les projets.
\end{itemize}

\section{Conclusion}

Cette mission m’a permis de participer à plusieurs projets numériques au sein de MEDIASSIVE, tout en occupant un rôle à la fois technique et organisationnel. À travers FinderrLink, ScanVivo et Batiscan, j’ai pu renforcer mes compétences en développement full-stack, en déploiement web et en coordination d’équipe. Cette expérience m’a également permis de mieux comprendre le fonctionnement d’un environnement professionnel orienté vers la création de solutions SaaS internes et la gestion de projets digitaux.
